\begin{center}
    {\LARGE \textbf{ĐỀ LUYỆN 04}}\\[6pt]
    {\large (Dành cho đội tuyển quốc gia)}
\end{center}

\begin{tabularx}{\textwidth}{|c|X|c|c|c|}
\hline
\textbf{Bài} & \textbf{Tên bài} & \textbf{File đầu vào} & \textbf{File đầu ra} & \textbf{Điểm} \\
\hline
1 & Chia và nối mảng & \texttt{SADD.INP} & \texttt{SADD.OUT} & 7 \\
\hline
2 & Rắn săn mồi & \texttt{SNAGR.INP} & \texttt{SNAGR.OUT} & 7 \\
\hline
3 & Ice-cream & \texttt{ICREAM.INP} & \texttt{ICREAM.OUT} & 6 \\
\hline
\end{tabularx}



\section*{Bài 1: Chia và nối mảng}
% https://www.codeThái.com/problems/SPLITADD

Bạn được cho ba số nguyên $N, L, R$. Biết rằng $N$ là một lũy thừa của $2$, tức là $N = 2^x$ với $1 \le x \le 60$.

Xét hàm $f$ tác động lên một mảng $A$ như sau:

\begin{itemize}
    \item Nếu $A$ chỉ có một phần tử, thì $f(A)=A$.
    \item Ngược lại, đặt
    \[
        H_1 = [A_1, A_3, A_5, \ldots, A_{N-1}], \qquad
        H_2 = [A_2, A_4, A_6, \ldots, A_N].
    \]
    Khi đó:
    \[
        f(A) = f(H_1) + f(H_2)
    \]
    trong đó dấu “+” là phép nối hai mảng.
\end{itemize}

Ban đầu, ta có mảng $A = [1, 2, 3, \ldots, N]$.  
Gọi $B = f(A)$.

Hãy tính:
\[
\sum_{i=L}^{R} B_i
\]
và in ra kết quả theo modulo $10^9+7$.

\subsection*{Dữ liệu vào}

Dòng đầu: số nguyên $T$ — số lượng test.

Mỗi test gồm một dòng chứa ba số nguyên $N, L, R$.

\subsection*{Dữ liệu ra}

Với mỗi test, in ra giá trị $\sum_{i=L}^{R} B_i \bmod (10^9+7)$.

\subsection*{Ràng buộc}

\begin{itemize}
    \item $1 \le T \le 10^5$
    \item $N = 2^x$
    \item $1 \le x \le 60$
    \item $1 \le L \le R \le N$
\end{itemize}

\subsection*{Ví dụ}

\textbf{Input}
\begin{verbatim}
4
8 7 7
8 1 7
8 3 6
1099511627776 1 100
\end{verbatim}

\textbf{Output}
\begin{verbatim}
4
28
18
270693356
\end{verbatim}

\subsection*{Giải thích}

\textbf{Test case 1.}

Với $N=8$:

\[
A = [1,2,3,4,5,6,7,8]
\]

Khi đó:

\[
f(A) = [1,5,3,7,2,6,4,8].
\]

\newpage

\section*{Bài 2. Rắn săn mồi}
% https://www.codeThái.com/problems/SNAGR

Trên một bảng kích thước $N \times N$, các hàng được đánh số từ $1$ đến $N$ từ trên xuống, các cột được đánh số từ $1$ đến $N$ từ trái sang. Ô ở hàng $i$, cột $j$ được ký hiệu là $(i,j)$.

Thái đứng tại ô $(1,1)$ và muốn đi đến ô $(N,N)$. Mỗi bước, Thái có thể di chuyển lên/xuống/trái/phải sang ô kề cạnh, miễn là không ra ngoài bảng.  
Tức là từ ô $(i,j)$ Thái có thể đi sang một trong bốn ô:
\[
(i-1,j),\ (i+1,j),\ (i,j-1),\ (i,j+1).
\]

Trên bảng có một số con rắn.  
Mỗi con rắn bắt đầu tại ô $(x_i,y_i)$ và có hướng $d_i \in \{U, D, L, R\}$ (lên, xuống, trái, phải). Cơ chế hoạt động như sau:

\begin{itemize}
    \item Ở thời điểm ban đầu, rắn chỉ chiếm đúng ô $(x_i,y_i)$.
    \item Cuối mỗi giây, nó mở rộng thêm một đơn vị theo hướng của mình.
\end{itemize}

Ví dụ: rắn tại $(4,7)$ có hướng \texttt{L} sẽ chiếm các ô:
\[
(4,7),\ (4,6),\ (4,5),\ldots
\]
cho đến khi chạm biên bảng.

Nhiều rắn có thể chiếm cùng một ô.

Thái gọi một ô là \textbf{nguy hiểm} nếu có ít nhất một rắn chiếm ô đó tại thời điểm đang xét.

Độ nguy hiểm của một đường đi từ $(1,1)$ đến $(N,N)$ là \textbf{số lượng ô nguy hiểm nằm trên đường đi đó}.

Bạn cần trả lời $Q$ truy vấn:
\[
\text{Với thời điểm } T,\ \text{hỏi độ nguy hiểm nhỏ nhất của một đường đi từ } (1,1) \text{ đến } (N,N) \text{ tại thời điểm } T?
\]

\subsection*{Dữ liệu vào}

\begin{itemize}
    \item Dòng đầu: hai số nguyên $N$ và $Q$.
    \item $N$ dòng tiếp theo: dòng thứ $i$ là xâu $S_i$ độ dài $N$, mỗi ký tự thuộc \texttt{.LRUD}.
    \item $Q$ dòng tiếp theo: dòng thứ $i$ chứa một số nguyên $T_i$.
\end{itemize}

\subsection*{Dữ liệu ra}

In ra $Q$ dòng, mỗi dòng một số nguyên — đáp án cho truy vấn tương ứng.

\subsection*{Ràng buộc}

\begin{itemize}
    \item $1 \le N \le 300$
    \item $1 \le Q \le 2 \cdot 10^5$
    \item $1 \le T_i \le 10^6$
    \item Mỗi $S_i$ chỉ chứa các ký tự \texttt{.LRUD}
\end{itemize}

\subsection*{Sample Input}

\begin{verbatim}
5 3
....L
..D..
.....
R..U.
.....
1
3
6
\end{verbatim}

\subsection*{Sample Output}

\begin{verbatim}
0
1
3
\end{verbatim}

\begin{figure}[htbp]
  \centering
  \includegraphics[scale=0.25]{De04_bai2_test1.png}
  \caption{Chú thích test 1}
  \label{fig:myimage}
\end{figure}

\newpage

\section*{Bài 3: Ice-cream}

Có $N$ cây kem đặt thành một hàng, trong đó cây kem thứ $i$ có giá tiền $C[i]$ và hương vị $F[i]$.

Thái muốn mua một \textbf{đoạn liên tiếp} các cây kem sao cho tổng giá tiền của đoạn đó không vượt quá ngân sách $K$.  
Sau đó, Thái có thể chọn \textbf{một tập con bất kỳ} (có thể rỗng, trừ khi yêu cầu) của các cây kem đã mua để trộn thành một cây kem mới.  

Nếu chọn $m$ cây kem có hương vị $A_1, A_2, \ldots, A_m$, thì hương vị khi trộn sẽ bằng:
\[
A_1 \oplus A_2 \oplus \cdots \oplus A_m,
\]
trong đó $\oplus$ là phép XOR theo bit.

Nhiệm vụ của bạn:  
\[
\textbf{Tìm giá trị hương vị lớn nhất mà Thái có thể đạt được.}
\]

Đảm bảo rằng có ít nhất một cây kem có giá không vượt quá $K$.

\subsection*{Dữ liệu vào}

\begin{itemize}
    \item Dòng đầu chứa số nguyên $T$ — số bộ test.
    \item Mỗi bộ test:
    \begin{itemize}
        \item Dòng đầu chứa hai số nguyên $N$ và $K$.
        \item $N$ dòng tiếp theo: dòng $i$ chứa hai số nguyên $C[i]$ và $F[i]$.
    \end{itemize}
\end{itemize}

\subsection*{Dữ liệu ra}

Với mỗi bộ test, in ra một số nguyên — hương vị tối đa Thái có thể đạt được.

\subsection*{Ràng buộc}

\begin{itemize}
    \item $1 \le T \le 10000$
    \item $1 \le N \le 100000$
    \item Tổng $N$ trong tất cả các test không vượt quá $100000$
    \item $1 \le K \le 10^{14}$
    \item $1 \le C[i] \le 10^9$
    \item $1 \le F[i] \le 10^9$
    \item Tồn tại ít nhất một $i$ sao cho $C[i] \le K$
\end{itemize}

\subsection*{Sample Input}

\begin{verbatim}
2
5 15
2 4
7 1
5 9
15 6
3 12
3 1000
123 8
99999 1
80 7
\end{verbatim}

\subsection*{Sample Output}

\begin{verbatim}
13
8
\end{verbatim}

\subsection*{Giải thích}

\textbf{Test 1.}  
Các hương vị dạng nhị phân:  
\[
(0100, 0001, 1001, 0110, 1100)
\]
giá tương ứng:
\[
(2, 7, 5, 15, 3)
\]

Thái có thể mua cây thứ $5$ và được hương vị $12$.  
Nhưng nếu mua 3 cây đầu (tổng giá $2+7+5 = 14 \le 15$), Thái có thể chọn trộn cây 1 và 3:
\[
0100 \oplus 1001 = 1101_2 = 13.
\]
Không thể đạt kết quả cao hơn.

\textbf{Test 2.}  
Thái không thể mua cây thứ 2 vì giá quá cao, nên chỉ có thể mua cây 1 hoặc cây 3.  
Mua cây 1 cho hương vị $8$, là phương án tối ưu.
